\subsection{Sprint 4}

A \textit{Sprint} 4 já estava prevista, desde o planejamento inicial do projeto, como uma etapa de consolidação. Com o sistema praticamente finalizado, a proposta era concentrar os esforços na correção de \textit{bugs} e, havendo margem, ampliar o escopo com funcionalidades adicionais. Para preparar essa fase, solicitei à equipe que utilizasse o aplicativo até a aula seguinte, avaliando a usabilidade e levantando eventuais novos problemas antes da distribuição das tarefas.

No escopo formal constavam apenas as atividades referentes à tela de perfil do usuário, tanto a visualização quanto a edição. Como entrega adicional, combinamos com as \textit{stakeholders} a tela de \textit{dashboard} do portal de administrador. Iniciamos a \textit{sprint} com quinze tarefas, das quais cinco correspondiam a \textit{bugs} já mapeados.

Mantivemos a abordagem que funcionara bem na \textit{sprint} anterior: revisamos as tarefas das novas \textit{features} entre os \ac{ages} IV e, em seguida, novamente com todo o time durante a \textit{planning}. Essa revisão coletiva levantou ajustes pertinentes, como o desmembramento da edição de perfil em duas tarefas, a fim de reduzir seu tamanho. Uma passou a tratar exclusivamente do \textit{patch} das preferências do usuário, e a outra, da edição da foto de perfil. Por se tratar de uma \textit{sprint} com forte componente de correção, optamos por dar liberdade à equipe na escolha das tarefas, permitindo que quem havia implementado uma funcionalidade com defeito pudesse assumir sua correção.

O engajamento do time ficou evidente desde o princípio, com tarefas sendo entregues ainda na primeira semana. Diante desse ritmo, em conjunto com os demais \ac{ages} IV, criamos novas \textit{features} com prioridade baixa, estabelecendo que só poderiam ser iniciadas após a conclusão das pendências existentes. Para compor uma entrega mais completa, desenvolvemos recuperação de senha, gestão de trilhas e gestão de usuários para o portal de administrador.

À medida que a equipe passou a utilizar mais o aplicativo, surgiram novos defeitos. Os cinco \textit{bugs} conhecidos no início da \textit{sprint} chegaram a doze em sua metade, todos ligados a pequenos detalhes. Ainda assim, a \textit{sprint} transcorreu de forma tranquila, e foi nela que percebi a maior dedicação do time ao longo de todo o projeto. As entregas se concentraram bem antes do \textit{code freeze}, e ao alcançá-lo já tínhamos todas as tarefas concluídas: entregamos a correção dos \textit{bugs} e as funcionalidades previstas no escopo e, ainda, as \textit{features} adicionais incorporadas como \textit{overdelivery}. Paralelamente, um dos \ac{ages} III trabalhou no \textit{deploy} via TestFlight, ferramenta de testes da Apple, o que permitiu disponibilizar o aplicativo às \textit{stakeholders} já na entrega e agregou valor significativo ao projeto.

Encerrado o \textit{code freeze}, reuni-me com os demais \ac{ages} IV para montar a apresentação final e executar um teste completo do fluxo do aplicativo e do portal web. Identificamos alguns defeitos, concentrados sobretudo no portal de administrador, sendo que um deles comprometia a importação de \acp{csv} de questões e simulados. Esse comportamento decorria de um \textit{bug} apontado ainda no fim da \textit{Sprint} 3, em que a importação esperava o \ac{uuid} das disciplinas e da trilha. O \textit{bug} já havia sido corrigido, porém apenas no \textit{backend}; no \textit{frontend}, o código ainda não havia sido alterado. Por estarmos nos últimos dias, assumi tanto essa correção quanto a montagem das listas de importação de questões e simulados.

A correção em si não apresentou dificuldades; o trabalho mais demorado foi a preparação das listas. Como as provas do ENEM são disponibilizadas em PDF de duas colunas, as ferramentas convencionais de extração não processavam o arquivo adequadamente. Encontrei, então, uma biblioteca em Python que atendia exatamente a essa necessidade: separava as questões em arquivos \ac{json}, com suas respectivas respostas corretas, e organizava as imagens em uma pasta dedicada. A partir dela, montei os \acp{csv} com as questões do ENEM de 2017 a 2021 para o banco de questões e utilizei a edição de 2022 para compor um simulado mais enxuto, com vinte e duas questões. As importações ocorreram sem falhas, e o projeto foi concluído.

Na segunda-feira, data prevista para a apresentação, confirmei a presença das clientes, e elas informaram que não poderiam comparecer em razão de provas finais. A complicação era que a quarta-feira, último dia disponível, coincidia com uma partida da Copa do Mundo, o que inviabilizava o horário regular da aula. Acordamos antecipar a apresentação para as 16h30, e elas concordaram. Aproveitamos essa última aula para tratar três defeitos recém-identificados: a alteração da foto de perfil não funcionava; tópicos excluídos por um administrador continuavam visíveis aos usuários; e as imagens dos tópicos não carregavam corretamente. Dois deles foram resolvidos ainda em aula, restando pendente apenas a edição da foto de perfil, corrigida a tempo da apresentação.

A apresentação ocorreu bem, ainda que com a presença de apenas uma das \textit{stakeholders}. Ela se mostrou bastante satisfeita com o resultado e visivelmente animada com tudo o que foi entregue além do combinado, demonstrando grande contentamento com o desfecho do projeto.
